\documentclass[11pt]{article}

\usepackage[margin=1in,headheight=22pt,headsep=20pt,footskip=30pt]{geometry}
\usepackage[T1]{fontenc}
\usepackage{lmodern}
\usepackage{xcolor}
\usepackage{hyperref}
\usepackage{fancyhdr}

\definecolor{NFPrimary}{HTML}{242632}
\definecolor{NFAccent}{HTML}{3DDCFF}
\definecolor{NFLink}{HTML}{007F95}
\definecolor{NFInk}{HTML}{242632}
\definecolor{NFMuted}{HTML}{656B78}
\definecolor{NFLight}{HTML}{F1F4F7}

\hypersetup{
  pdftitle={Neofluxon User Guide},
  pdfauthor={Iurie Nistor},
  colorlinks=true,
  linkcolor=NFLink,
  urlcolor=NFLink
}

\title{Neofluxon User Guide}
\author{Iurie Nistor}
\date{}

\pagestyle{fancy}
\fancyhf{}
\fancyhead[L]{\small\sffamily\bfseries\color{NFPrimary}NEOFLUXON}
\fancyhead[R]{\small\sffamily\color{NFMuted}USER GUIDE}
\fancyfoot[C]{\small\sffamily\color{NFMuted}\thepage}
\renewcommand{\headrulewidth}{0.4pt}
\renewcommand{\footrulewidth}{0pt}
\renewcommand{\headrule}{\hbox to\headwidth{\color{NFAccent}\leaders\hrule height \headrulewidth\hfill}}

\setlength{\leftmargini}{1.6em}
\setlength{\leftmarginii}{1.8em}

\begin{document}

\hypersetup{pageanchor=false}
\makeatletter
\begin{titlepage}
\thispagestyle{empty}
\color{NFInk}
\vspace*{1.2cm}
\noindent\textcolor{NFAccent}{\rule{3.2cm}{3pt}}

\vspace{1.4cm}
\begin{flushleft}
{\sffamily\bfseries\fontsize{32}{38}\selectfont\color{NFPrimary}\@title\par}
\vspace{0.5cm}
{\sffamily\Large\color{NFMuted}Author: \@author\par}
\end{flushleft}

\vspace{1.3cm}
\noindent\colorbox{NFLight}{%
  \parbox{\dimexpr\textwidth-2\fboxsep\relax}{%
    \vspace{0.35cm}
    \textbf{Software version:} @NEOFLUXON_DOC_VERSION@ \par
    \vspace{0.25cm}
    \textbf{License:} This work, including all images it uses, is released under
    the CC0 1.0 Universal (CC0 1.0) Public Domain Dedication. \par
    \vspace{0.25cm}
    \textbf{Website:} \url{https://neofluxon.com}
    \vspace{0.35cm}
  }}
\vfill
\noindent\textcolor{NFAccent}{\rule{\textwidth}{1pt}}
\end{titlepage}
\makeatother

\hypersetup{pageanchor=true}
\pagenumbering{roman}
\setcounter{tocdepth}{2}
\setcounter{secnumdepth}{0}
\tableofcontents
\clearpage
\pagenumbering{arabic}

\section{Introduction}

\subsection{What Is Neofluxon?}

Neofluxon is a \href{https://www.gnu.org/philosophy/free-sw.html}{Free
Software} desktop application for browsing photos in folders and organizing
photo libraries. It is licensed under the
\href{https://www.gnu.org/licenses/gpl-3.0.html}{GNU General Public License,
version 3 or later} and is available for
\href{https://www.gnu.org/gnu/linux-and-gnu.html}{GNU/Linux} and Windows.

Neofluxon currently supports RAW, JPEG, and PNG image files. It does not
modify original photos or change the existing folder structure on the
filesystem.

\subsection{How Neofluxon Works}

Neofluxon provides two ways to browse photos:
\begin{itemize}
  \item \textbf{Folder Mode} displays photos directly from their existing
        filesystem folders. It does not require a library or an import.
  \item \textbf{Library Mode} displays photos recorded in a library. A
        library stores catalog information about folders and photos; it
        does not copy or move the original files.
\end{itemize}

\section{Working with Neofluxon}

\subsection{The Main Window}

The main window is organized into three main areas:
\begin{itemize}
  \item \textbf{Left panel:} provides navigation. Its contents depend on
        whether Folder Mode or Library Mode is active.
  \item \textbf{Central browser:} displays photos from the current folder or
        selected library location.
  \item \textbf{Right panel:} displays metadata for the selected photo.
\end{itemize}

The left and right panels can be resized or collapsed.

\subsection{Switching Between Folder Mode and Library Mode}

Select \textbf{Folders} or \textbf{Library} from the Main Menu in the left panel to switch between browsing modes.

In Folder Mode, select a folder in the source panel to display its photos in the central browser.

In Library Mode, select a library and then choose one of its representations to browse its contents.
The currently selected library is always displayed in the top bar.

\subsection{Keyboard Navigation}

When the central browser is focused, use the arrow keys to navigate photos:
\textbf{Left} and \textbf{Right} move to the previous or next photo, while
\textbf{Up} and \textbf{Down} move between rows in Grid mode. Press
\textbf{Home} or \textbf{End} to move to the first or last thumbnail, or
\textbf{Page Up} and \textbf{Page Down} to move by a page in Grid mode.

\subsection{Selecting and Viewing Photos}

Select a photo thumbnail in the central browser to display its metadata in the right
panel. Double-click a thumbnail to open browser Preview Mode. In Preview Mode, select
another thumbnail in the filmstrip to view a different photo.

\section{Browsing Photos in Folders}

\subsection{Finding a Folder}

The Folder Mode source panel provides three ways to locate a directory.

\subsubsection{Folder Tree}

Expand the folder tree to navigate the filesystem. Double-click a folder to
show its contents in the central browser.

\subsubsection{Places}

\textbf{Places} is a list of standard locations provided by the operating
system. Selecting a place opens it in the folder tree.

\subsubsection{Storage}

\textbf{Storage} is list of available drives or volumes. Selecting a
storage location opens it in the folder tree. Currently, there is need to restart
the applicaiton for changes to connected storage devices
to be reflected in the list.

\subsection{Browsing a Folder}

Select a folder in the source panel to show its photos in the
central browser. The enabled file types in Settings determine which files
are show.

\subsection{Recursive and Non-Recursive Browsing}

With recursive scanning disabled, Neofluxon lists photos in the selected
folder only. When recursive-scan is enabled the browser will include
photos in subfolders.

The recursive scan control is located in the browser's
top bar and is available only in Folder Mode.

\section{Working with Libraries}

Neofluxon supports multiple libraries in the same application instance.
The user can switch between libraries at runtime, but only one library
can be selected at a time. The currently selected library is
shown in the application's top bar.

A library consists of three representations: \textbf{Datetime}, \textbf{Folders}, and \textbf{Equipment}.
It keeps a catalog of folders and photos so they can be browsed
using different representations.

Importing a folder adds its supported photos to the selected library.
Neofluxon does not copy or move the source files.

\subsection{Creating and Selecting a Library}

The library list appears in the left panel when Library Mode is active.
Select a library to make it current.

Use the plus button beside \textbf{Libraries} to create a library and enter
its name.

\subsection{Importing Photo Folders}

To add a folder to a library, select the library, switch to
\textbf{Folders} in the main menu, and locate the folder. Select the plus button displayed
when the pointer is over its row. The folder's supported photos are added to
the library catalog; the original files remain in their existing locations.

The status bar displays background activity, including folder imports.
During an import, it reports the numbers of discovered, imported, and failed
photos. Completed are removed from the active-task list.

\subsection{Library Representations}

After selecting a library, choose a representation to browse its contents.

\subsubsection{Date and Time}

This representation groups photos by the \textbf{date taken} recorded in their metadata.

\subsubsection{Folder}

This representation shows the folders and subfolders added to the library.

\subsubsection{Equipment}

This representation groups photos by camera or lens information when that
information is available.

\subsection{Navigating Library Contents}

The library contents tree shows entries for the selected representation.
Select a date, folder, camera, or lens entry to display its photos in the
central browser.

\subsection{Managing Libraries}

Right-click a library to rename or delete it. Deleting a library removes its
catalog entries and folders that are not shared with another library.
It does not delete the source photo files or folders on the filesystem.

\section{Browsing and Viewing Photos}

The central browser displays photos from the current folder or selected
library location. It provides Grid Mode and Preview Mode, along with controls
for sorting and filtering.

\subsection{Grid Mode}

Grid Mode displays photo thumbnails. Select a thumbnail to select a photo;
double-click it to open Preview Mode.

\subsection{Preview Mode and Filmstrip}

Preview Mode displays a larger view of the selected photo and a filmstrip
containing other photos. Select another thumbnail in the filmstrip to change
the preview. Use the Grid view control to return to Grid Mode.

\subsection{Sorting Photos}

Use the sort menu to choose one of the available sort fields:
\begin{itemize}
  \item Name
  \item Date Taken
  \item File Size
\end{itemize}

Use the arrow control to change the sort direction.

\subsection{Filtering Photos}

Use the filter menu to show selected photo formats. RAW, JPEG, and PNG
formats can be selected, including multiple formats at the same time. Choose
\textbf{Clear format filters} to remove the format filters.

\section{Viewing Photo Information}

Select a photo in the central browser to display its available metadata in
the right panel.

\subsection{Metadata Categories}

The metadata panel contains collapsible sections:
\begin{itemize}
  \item \textbf{Image:} name, path, size, format, and dimensions.
  \item \textbf{Capture:} date taken, camera, lens, and exposure values.
  \item \textbf{Location:} GPS information, when present. Double-clicking a
        location opens the system file manager and selects the photo.
  \item \textbf{Organization:} rating, tags, collections, and project.
  \item \textbf{Description:} title, description, author, copyright, and
        keywords.
  \item \textbf{Technical:} bit depth, color space, orientation, and other
        file properties.
\end{itemize}

The information shown depends on the metadata available in each photo.

\section{Application Settings}

Open Settings using the button at the right side of the application top bar.

\subsection{File Types}

The \textbf{File Types} page controls which supported image formats are
enabled. Enabled formats determine which photos are included when browsing
folders or importing them into a library.

\subsection{Performance}

The \textbf{Performance} page contains controls for thumbnail and preview
memory limits, the disk-cache location and size, and the minimum free space
to keep on the cache drive.

\section{Appendix}

\subsection{Supported Image Formats}

Neofluxon recognizes the following file extensions:
\begin{itemize}
  \item \textbf{RAW:} Canon CR2 (\texttt{.cr2}) and CR3 (\texttt{.cr3}),
        Adobe DNG (\texttt{.dng}), Nikon NEF (\texttt{.nef}), and Fujifilm
        RAF (\texttt{.raf}).
  \item \textbf{Raster:} JPEG (\texttt{.jpg}, \texttt{.jpeg}) and PNG
        (\texttt{.png}).
\end{itemize}

\subsection{Cache Overview}

Neofluxon uses caches to retain thumbnail and preview data for reuse:
\begin{itemize}
  \item \textbf{Display cache} stores ready-to-display thumbnail pixmaps so
        the browser can render them.
  \item \textbf{Memory caches} hold decoded thumbnail and preview images in
        memory while the application is running.
  \item \textbf{Disk cache} stores regenerable thumbnail and preview data
        between application sessions.
\end{itemize}

\subsection{Building and Installing from Source}

To build Neofluxon from source, install CMake 3.31 or newer, a C++26
compiler, Qt 6, and the development dependencies listed in the repository
README. Configure, build, and install with:
\begin{verbatim}
cmake -S . -B build -DCMAKE_BUILD_TYPE=Release
cmake --build build --parallel
cmake --install build
\end{verbatim}

To build the guide PDF as part of the application build, enable documentation
generation during configuration:
\begin{verbatim}
cmake -S . -B build -DNEOFLUXON_DOCUMENTATION=ON
\end{verbatim}

Documentation generation requires \texttt{pdflatex}. The resulting PDF is
installed in the platform's documentation directory.

\end{document}
